%========================================
% MatinBook — Stage 08: Code System Test
% Version: v1.1
% Purpose:
%   Validate the code display system of MatinBook v1.1:
%     1. minted integration with Pygments
%     2. Python, C++, and plain-text listings
%     3. Inline code (\inlcode)
%     4. Code captions and numbering (\codecaption)
%     5. Cross-references to code (\coderef, \cref)
%     6. Persian comments inside code (\pc)
%     7. Code font (JetBrains Mono with ligatures)
%
% Compilation:
%   xelatex -shell-escape stage08-code.tex
%   xelatex -shell-escape stage08-code.tex
%   (2 passes required for cover + cross-references)
%
% Expected outcome:
%   A professionally typeset PDF demonstrating that code
%   listings render with syntax highlighting, correct font,
%   proper captions, and working cross-references.
%========================================

%------------------------------------------------------------------------
% Search Path (for tests directory)
%------------------------------------------------------------------------
\makeatletter
\def\input@path{{../}}
\makeatother

%------------------------------------------------------------------------
% Document Class
%------------------------------------------------------------------------
\documentclass{matinbook}

%------------------------------------------------------------------------
% Book Metadata
%------------------------------------------------------------------------
\title{آزمون سیستم نمایش کد}
\author{تیم توسعه MatinBook}
\date{مهر ۱۴۰۵}

\booktitle{آزمون سیستم نمایش کد}

\renewcommand{\repository}{github.com/matinmahmoudi-cs/matinbook}

%========================================================================
% DOCUMENT
%========================================================================
\begin{document}

%------------------------------------------------------------------------
% FRONT COVER
%------------------------------------------------------------------------
\makecover
    {آزمون سیستم نمایش کد}
    {ارزیابی رنگ‌آمیزی، کپشن و ارجاع در MatinBook}
    {تیم توسعه MatinBook}
    {مهر ۱۴۰۵}

%------------------------------------------------------------------------
% FRONT MATTER (symmetric margins)
%------------------------------------------------------------------------
\frontmatter
\frontmattergeometry

% Title page
\maketitle

% Copyright / Colophon
\thispagestyle{empty}
\vfill
\begin{center}
    {\large\textbf{شناسنامه آزمون}}

    \vspace{0.8cm}

    \begin{tabular}{rl}
        عنوان: & آزمون سیستم نمایش کد \\
        نسخه: & \matinbookversion \\
        تاریخ: & \matinbookdate \\
        موتور: & \lr{XeLaTeX} \\
        مجوز: & \lr{MIT} \\
    \end{tabular}

    \vspace{0.8cm}

    {\small این سند بخشی از مجموعه آزمون‌های یکپارچگی \lr{MatinBook} است.}
\end{center}
\clearpage

% Preface
\chapter*{پیشگفتار}
\addcontentsline{toc}{chapter}{پیشگفتار}

این سند، هشتمین آزمون از مجموعه‌ی آزمون‌های یکپارچگی
\lr{MatinBook} نسخه‌ی \lr{\matinbookversion} است. هدف آن،
اعتبارسنجی کامل سیستم نمایش کد کلاس در هفت حوزه‌ی متمایز
است: یکپارچگی \lr{minted} با \lr{Pygments}، کدهای
\lr{Python} و \lr{C++}، کد درون‌خطی، کپشن و شماره‌گذاری
کد، ارجاع به کد، کامنت فارسی داخل کد، و فونت کد.

در کتاب‌های برنامه‌نویسی، نمایش حرفه‌ای کد اهمیت
ویژه‌ای دارد. کد باید با رنگ‌آمیزی نحوی مناسب، فونت
خوانا، شماره‌گذاری منظم و کپشن‌های گویا نمایش داده
شود. \lr{MatinBook} با استفاده از \lr{minted} و
فونت \lr{JetBrains Mono}، این اهداف را محقق می‌سازد.

هر بخش از این آزمون، با مثال‌های عملی و ارجاعات متقابل
همراه است تا خواننده بتواند درستی رفتار هر زیرسیستم را
در بافت واقعی یک کتاب ارزیابی کند.

\clearpage

% Table of Contents
\tableofcontents
\clearpage

% List of Figures
\listoffigures
\clearpage

% List of Tables
\listoftables
\clearpage

%------------------------------------------------------------------------
% MAIN MATTER (wide margin for notes)
%------------------------------------------------------------------------
\mainmatter
\mainmattergeometry

%========================================================================
% CHAPTER 1 — PYTHON PROGRAMMING
%========================================================================
\chapter{برنامه‌نویسی پایتون}
\label{chap:python}

\section{مبانی زبان}
\label{sec:python-basics}

پایتون\index{پایتون} یک زبان برنامه‌نویسی سطح بالا،
تفسیری و همه‌منظوره است که بر خوانایی کد تأکید دارد.
در ادامه، مفاهیم پایه را با کد بررسی می‌کنیم.

\mnote{پایتون در سال ۱۹۹۱ توسط \lr{Guido van Rossum} طراحی شد.}

\subsection{متغیرها و انواع داده}
\label{sec:python-variables}

در \cref{code:variables} با انواع داده‌ی پایه در پایتون
آشنا می‌شویم:

\begin{latin}
\begin{minted}{python}
# Basic data types in Python
age: int = 25                          # Integer type
temperature: float = 36.5              # Floating-point type
name: str = "Ali"                      # String type
is_active: bool = True                 # Boolean type

# Compound data structures
scores: list = [18, 17, 20, 19]       # List - ordered, mutable
coordinates: tuple = (35.7, 51.4)      # Tuple - ordered, immutable
student: dict = {                       # Dictionary - key-value pairs
    "name": "Ali",
    "age": 20,
    "grades": [18, 17, 20]
}

# Set - unordered collection of unique elements
unique_ids: set = {1001, 1002, 1003}
\end{minted}
\end{latin}
\codecaption{انواع داده در پایتون}
\label{code:variables}

\subsection{ساختارهای شرطی}
\label{sec:python-conditions}

استفاده از دستور \inlcode{if-elif-else} در پایتون در
\cref{code:conditions} نمایش داده شده است:

\begin{latin}
\begin{minted}{python}
def get_grade(score: float) -> str:
    """
    Convert numerical score to letter grade.

    Args:
        score: Numerical score between 0 and 100.

    Returns:
        Letter grade (A, B, C, D, or F).
    """
    if score >= 90:
        return 'A'
    elif score >= 80:
        return 'B'
    elif score >= 70:
        return 'C'
    elif score >= 60:
        return 'D'
    else:
        return 'F'


# Create a list of students with their scores
students = [
    {"name": "Ali", "score": 85},
    {"name": "Sara", "score": 92},
    {"name": "Reza", "score": 78},
]

# Calculate and display grades for each student
for student in students:
    grade = get_grade(student["score"])
    print(f"{student['name']}: {student['score']} -> {grade}")
\end{minted}
\end{latin}
\codecaption{ساختار شرطی و حلقه در پایتون}
\label{code:conditions}

\subsection{توابع و بازگشت}
\label{sec:python-recursion}

توابع بازگشتی در \cref{code:recursion} نشان داده شده‌اند:

\begin{latin}
\begin{minted}{python}
def factorial(n: int) -> int:
    """
    Calculate the factorial of n using recursion.

    Args:
        n: A non-negative integer.

    Returns:
        The factorial of n (n!).
    """
    # Base case: 0! = 1 and 1! = 1
    if n <= 1:
        return 1
    # Recursive case: n! = n * (n-1)!
    return n * factorial(n - 1)


def fibonacci(n: int) -> int:
    """
    Calculate the nth Fibonacci number using recursion.

    Args:
        n: The index (0-based) in the Fibonacci sequence.

    Returns:
        The nth Fibonacci number.
    """
    # Base cases
    if n <= 0:
        return 0
    elif n == 1:
        return 1
    # Recursive case: F(n) = F(n-1) + F(n-2)
    return fibonacci(n - 1) + fibonacci(n - 2)


# Display the first 10 Fibonacci numbers
for i in range(10):
    print(f"F({i}) = {fibonacci(i)}")
\end{minted}
\end{latin}
\codecaption{توابع بازگشتی در پایتون}
\label{code:recursion}

همان‌طور که در \cref{code:recursion} می‌بینیم، پیچیدگی
زمانی تابع فیبوناچی بازگشتی \lr{$O(2^{n})$} است که برای
$n$های بزرگ کارآمد نیست.

\subsection{برنامه‌نویسی شیءگرا}
\label{sec:python-oop}

برنامه‌نویسی شیءگرا در \cref{code:oop} با یک مثال
کامل نشان داده شده است:

\begin{latin}
\begin{minted}{python}
from dataclasses import dataclass
from typing import List


@dataclass
class Student:
    """
    A class representing a university student.

    Attributes:
        name: Full name of the student.
        student_id: Unique student identification number.
        grades: List of numerical grades (0-20 scale).
    """
    name: str
    student_id: int
    grades: List[float]

    @property
    def average(self) -> float:
        """Calculate the arithmetic mean of all grades."""
        if not self.grades:
            return 0.0
        return sum(self.grades) / len(self.grades)

    @property
    def letter_grade(self) -> str:
        """Convert numerical average to letter grade."""
        avg = self.average
        if avg >= 18:
            return 'A'
        elif avg >= 15:
            return 'B'
        elif avg >= 12:
            return 'C'
        elif avg >= 10:
            return 'D'
        else:
            return 'F'

    def add_grade(self, grade: float) -> None:
        """
        Add a new grade to the student's record.

        Args:
            grade: A numerical grade between 0 and 20.

        Raises:
            ValueError: If grade is outside the valid range.
        """
        if 0 <= grade <= 20:
            self.grades.append(grade)
        else:
            raise ValueError(
                f"Grade must be between 0 and 20, got {grade}"
            )


# Create student instances and add grades
student = Student("Ali Alizadeh", 1001, [18, 17, 19])
student.add_grade(20)

# Display student information
print(f"Student: {student.name}")
print(f"ID: {student.student_id}")
print(f"Grades: {student.grades}")
print(f"Average: {student.average:.2f}")
print(f"Letter Grade: {student.letter_grade}")
\end{minted}
\end{latin}
\codecaption{کلاس Student با DataClass و Property}
\label{code:oop}

\section{کد درون‌خطی}
\label{sec:python-inline}

می‌توانیم در متن به توابع و متغیرها اشاره کنیم:
تابع \inlcode{print()} برای چاپ خروجی، تابع \inlcode{len()}
برای محاسبه‌ی طول، و \inlcode{range()} برای تولید
دنباله‌ی اعداد.

همچنین می‌توان به ماژول‌های استاندارد مانند
\inlcode{collections}، \inlcode{itertools} و
\inlcode{functools} اشاره کرد.

%========================================================================
% CHAPTER 2 — C++ PROGRAMMING
%========================================================================
\chapter{برنامه‌نویسی \lr{C++}}
\label{chap:cpp}

\section{مفاهیم پیشرفته}
\label{sec:cpp-advanced}

\subsection{اشاره‌گرهای هوشمند}
\label{sec:cpp-smart-pointers}

اشاره‌گرهای هوشمند و \lr{RAII} در \cref{code:smart-pointers}
نمایش داده شده‌اند:

\begin{latin}
\begin{minted}{cpp}
#include <iostream>
#include <memory>
#include <string>

/**
 * A simple resource class that logs creation and destruction.
 * Demonstrates RAII (Resource Acquisition Is Initialization).
 */
class Resource {
private:
    std::string name;

public:
    explicit Resource(const std::string& n) : name(n) {
        std::cout << "Resource '" << name << "' created\n";
    }

    ~Resource() {
        std::cout << "Resource '" << name << "' destroyed\n";
    }

    void use() const {
        std::cout << "Using resource: " << name << "\n";
    }
};

int main() {
    std::cout << "=== Unique Pointer Demo ===\n";

    // unique_ptr: exclusive ownership, cannot be copied
    auto dbConnection = std::make_unique<Resource>("Database");
    dbConnection->use();

    // Transfer ownership using std::move
    auto dbConnection2 = std::move(dbConnection);
    if (!dbConnection) {
        std::cout << "Original pointer is now null\n";
    }
    dbConnection2->use();

    std::cout << "\n=== Shared Pointer Demo ===\n";

    // shared_ptr: shared ownership with reference counting
    auto cache = std::make_shared<Resource>("Cache");
    std::cout << "Reference count: " << cache.use_count() << "\n";

    {
        auto cacheCopy = cache;
        std::cout << "Reference count: " << cache.use_count() << "\n";
        cacheCopy->use();
    }

    std::cout << "Reference count: " << cache.use_count() << "\n";
    cache->use();

    std::cout << "\n=== End of main ===\n";
    return 0;
}
\end{minted}
\end{latin}
\codecaption{اشاره‌گرهای هوشمند و \lr{RAII} در \lr{C++}}
\label{code:smart-pointers}

\subsection{الگوها (\lr{Templates})}
\label{sec:cpp-templates}

الگوها و مفاهیم \lr{C++20} در \cref{code:templates}
نمایش داده شده‌اند:

\begin{latin}
\begin{minted}{cpp}
#include <iostream>
#include <vector>
#include <stdexcept>
#include <string>
#include <concepts>

/**
 * Concept: a type must be comparable using < and > operators.
 */
template<typename T>
concept Comparable = requires(T a, T b) {
    { a < b } -> std::convertible_to<bool>;
    { a > b } -> std::convertible_to<bool>;
    { a == b } -> std::convertible_to<bool>;
};

/**
 * A generic Stack data structure implemented using std::vector.
 */
template<typename T>
requires Comparable<T>
class Stack {
private:
    std::vector<T> data;

public:
    void push(const T& value) {
        data.push_back(value);
    }

    T pop() {
        if (data.empty()) {
            throw std::runtime_error("Stack underflow");
        }
        T value = data.back();
        data.pop_back();
        return value;
    }

    const T& top() const {
        if (data.empty()) {
            throw std::runtime_error("Stack is empty");
        }
        return data.back();
    }

    bool empty() const { return data.empty(); }
    size_t size() const { return data.size(); }
};

int main() {
    Stack<int> intStack;
    intStack.push(10);
    intStack.push(20);
    intStack.push(30);

    std::cout << "Stack size: " << intStack.size() << "\n";

    std::cout << "Popping elements: ";
    while (!intStack.empty()) {
        std::cout << intStack.pop() << " ";
    }
    std::cout << "\n";

    return 0;
}
\end{minted}
\end{latin}
\codecaption{الگوها (\lr{Templates}) و مفاهیم (\lr{Concepts}) در \lr{C++20}}
\label{code:templates}

%========================================================================
% CHAPTER 3 — ALGORITHM IMPLEMENTATION
%========================================================================
\chapter{پیاده‌سازی الگوریتم‌ها}
\label{chap:algorithms}

\section{الگوریتم‌های جستجو}
\label{sec:algorithms-search}

\subsection{جستجوی دودویی}
\label{sec:algorithms-binary}

جستجوی دودویی در \cref{code:binary-search} با مستندات
کامل پیاده‌سازی شده است:

\begin{latin}
\begin{minted}{python}
from typing import List, Optional


def binary_search(arr: List[int], target: int) -> Optional[int]:
    """
    Perform binary search on a sorted array.

    Args:
        arr: A sorted list of integers (ascending order).
        target: The value to search for.

    Returns:
        The index of the target if found, None otherwise.

    Time Complexity: O(log n)
    Space Complexity: O(1)
    """
    left, right = 0, len(arr) - 1

    while left <= right:
        mid = left + (right - left) // 2

        if arr[mid] == target:
            return mid
        elif arr[mid] < target:
            left = mid + 1
        else:
            right = mid - 1

    return None


if __name__ == "__main__":
    sorted_array = [1, 3, 5, 7, 9, 11, 13, 15, 17, 19]
    test_cases = [7, 1, 19, 8]

    for target in test_cases:
        result = binary_search(sorted_array, target)
        if result is not None:
            print(f"Target {target} found at index {result}")
        else:
            print(f"Target {target} not found")
\end{minted}
\end{latin}
\codecaption{پیاده‌سازی جستجوی دودویی با مستندات کامل}
\label{code:binary-search}

\subsection{مرتب‌سازی سریع}
\label{sec:algorithms-quicksort}

مرتب‌سازی سریع در \cref{code:quicksort} پیاده‌سازی شده است:

\begin{latin}
\begin{minted}{python}
from typing import List


def quicksort(arr: List[int]) -> List[int]:
    """
    Sort a list using the quicksort algorithm.

    Args:
        arr: A list of integers to be sorted.

    Returns:
        A new sorted list in ascending order.

    Time Complexity: O(n log n) average, O(n^2) worst case
    Space Complexity: O(n) for the recursive call stack
    """
    if len(arr) <= 1:
        return arr

    pivot = arr[len(arr) // 2]
    left = [x for x in arr if x < pivot]
    middle = [x for x in arr if x == pivot]
    right = [x for x in arr if x > pivot]

    return quicksort(left) + middle + quicksort(right)


if __name__ == "__main__":
    unsorted = [64, 34, 25, 12, 22, 11, 90, 5]
    print(f"Original array: {unsorted}")

    sorted_arr = quicksort(unsorted)
    print(f"Sorted array:   {sorted_arr}")

    expected = sorted(unsorted)
    assert sorted_arr == expected, "Sorting failed!"
    print("Verification passed")
\end{minted}
\end{latin}
\codecaption{پیاده‌سازی مرتب‌سازی سریع با تأیید صحت}
\label{code:quicksort}

%========================================================================
% CHAPTER 4 — CODE DISPLAY FEATURES
%========================================================================
\chapter{ویژگی‌های نمایش کد}
\label{chap:code-features}

\section{تنظیمات نمایش}
\label{sec:code-settings}

\lr{MatinBook} از بسته‌ی \lr{minted} با تنظیمات زیر برای
نمایش کد استفاده می‌کند:

\begin{itemize}
    \item \textbf{فونت:} \lr{JetBrains Mono} با اندازه‌ی \lr{footnotesize}
    \item \textbf{شماره خطوط:} فعال (\lr{linenos=true})
    \item \textbf{شکستن خطوط:} فعال (\lr{breaklines=true})
    \item \textbf{کادر:} خطوط بالا و پایین (\lr{frame=lines})
    \item \textbf{رنگ کادر:} \lr{matinbordergray}
    \item \textbf{پس‌زمینه:} \lr{matinlightgray!30}
    \item \textbf{اندازه Tab:} ۴ فاصله (\lr{tabsize=4})
    \item \textbf{حذف تورفتگی مشترک:} فعال (\lr{autogobble=true})
    \item \textbf{استایل رنگ‌آمیزی:} \lr{friendly}
\end{itemize}

\mnote{کدهای داخل \lr{minted} باید همیشه داخل محیط \lr{latin} قرار گیرند.}

\section{ارجاع به کدها}
\label{sec:code-refs}

قابلیت ارجاع به قطعات کد با استفاده از \lr{\textbackslash coderef}:

\begin{itemize}
    \item \cref{code:variables} — انواع داده در پایتون
    \item \cref{code:conditions} — ساختارهای شرطی
    \item \cref{code:recursion} — توابع بازگشتی
    \item \cref{code:oop} — برنامه‌نویسی شیءگرا
    \item \cref{code:smart-pointers} — اشاره‌گرهای هوشمند \lr{C++}
    \item \cref{code:templates} — الگوها و مفاهیم \lr{C++20}
    \item \cref{code:binary-search} — جستجوی دودویی
    \item \cref{code:quicksort} — مرتب‌سازی سریع
\end{itemize}

\section{کامنت فارسی در کد}
\label{sec:code-persian-comments}

با استفاده از دستور \lr{\textbackslash pc}، می‌توان
کامنت فارسی داخل کد نوشت:

\begin{latin}
\begin{minted}{python}
# |\pc{این یک کامنت فارسی است}|
x = 5
y = 10
result = x + y  # |\pc{جمع دو متغیر}|
print(result)
\end{minted}
\end{latin}
\codecaption{کامنت فارسی داخل کد پایتون}
\label{code:persian-comments}

\section{مقایسه زبان‌ها}
\label{sec:code-comparison}

\begin{latin}
\begin{minted}{text}
Language     | Typing    | Paradigm            | Speed    | Use Case
-------------|-----------|---------------------|----------|------------------
Python       | Dynamic   | Multi-paradigm      | Medium   | Data Science, AI
C++          | Static    | Multi-paradigm      | Fast     | Systems, Games
Rust         | Static    | Functional/Imp      | Fast     | Systems, Web
JavaScript   | Dynamic   | Event-driven        | Medium   | Web Development
Go           | Static    | Concurrent          | Fast     | Cloud, Network
\end{minted}
\end{latin}
\codecaption{مقایسه زبان‌های برنامه‌نویسی}
\label{code:comparison}

%========================================================================
% CHAPTER 5 — SUMMARY
%========================================================================
\chapter{خلاصه‌ی نتایج}
\label{chap:summary}

\section{وضعیت سیستم کد}
\label{sec:summary-code-status}

جدول \cref{tab:code-status} وضعیت قابلیت‌های نمایش کد
\lr{MatinBook} را نشان می‌دهد.

\begin{matintable}{وضعیت قابلیت‌های نمایش کد}{tab:code-status}
\begin{tblr}{
    width = \textwidth,
    colspec = {l X[c] X[c]},
    hlines, vlines,
    colsep = 8pt,
    rowsep = 4pt,
    row{1} = {font=\bfseries},
}
قابلیت & توضیحات & وضعیت \\
رنگ‌آمیزی نحوی & پشتیبانی از ۳۰۰+ زبان برنامه‌نویسی & \textcolor{matingreen}{فعال} \\
شماره خطوط & نمایش شماره کنار هر خط کد & \textcolor{matingreen}{فعال} \\
کادر & خطوط بالا و پایین با رنگ \lr{matinbordergray} & \textcolor{matingreen}{فعال} \\
شکستن خطوط & خطوط طولانی خودکار می‌شکنند & \textcolor{matingreen}{فعال} \\
کد درون‌خطی & \lr{\textbackslash inlcode} برای ارجاع در متن & \textcolor{matingreen}{فعال} \\
ارجاع به کد & \lr{\textbackslash cref} با شماره‌ی خودکار & \textcolor{matingreen}{فعال} \\
کامنت فارسی & \lr{\textbackslash pc} داخل \lr{minted} & \textcolor{matingreen}{فعال} \\
فونت کد & \lr{JetBrains Mono} با لیگچرها & \textcolor{matingreen}{فعال} \\
\end{tblr}
\end{matintable}

\section{معیارهای ارزیابی}
\label{sec:summary-criteria}

در این آزمون، سیستم نمایش کد \lr{MatinBook} بر اساس
پنج معیار ارزیابی شد:

\begin{enumerate}
    \item \textbf{رنگ‌آمیزی}: دقت رنگ‌ها برای هر زبان.
    \item \textbf{خوانایی}: فونت و فاصله‌گذاری مناسب.
    \item \textbf{کپشن}: شماره‌گذاری منظم و ارجاع هوشمند.
    \item \textbf{چندزبانی}: پشتیبانی از زبان‌های مختلف.
    \item \textbf{کامنت فارسی}: یکپارچگی با متن فارسی.
\end{enumerate}

\section{نتیجه‌گیری}
\label{sec:summary-conclusion}

\textbf{تمام زیرسیستم‌های نمایش کد \lr{MatinBook} نسخه‌ی
\lr{\matinbookversion} با موفقیت بارگذاری و آزمایش شدند.
کدهای \lr{Python}، \lr{C++} و متن، با رنگ‌آمیزی مناسب،
مستندات کامل، کپشن‌های گویا و قابلیت ارجاع به‌درستی
کار می‌کنند.}

%------------------------------------------------------------------------
% BACK MATTER (symmetric margins)
%------------------------------------------------------------------------
\backmatter
\frontmattergeometry

% Bibliography (empty placeholder)
\printbibliography[title={منابع و مراجع}]

% Index
\printindex

%------------------------------------------------------------------------
% BACK COVER
%------------------------------------------------------------------------
\authorbio{%
    تیم توسعه‌ی \lr{MatinBook}، گروهی از پژوهشگران و برنامه‌نویسان
    فارسی‌زبان است که با هدف ارائه‌ی یک راه‌حل حرفه‌ای برای
    حروف‌چینی کتاب‌های فنی فارسی فعالیت می‌کند. این پروژه حاصل
    سال‌ها تجربه در حوزه‌ی نشر دیجیتال و برنامه‌نویسی است.
}

\makebackcover{%
    این سند، هشتمین آزمون از مجموعه‌ی آزمون‌های یکپارچگی
    \lr{MatinBook} نسخه‌ی \lr{\matinbookversion} است.

    \vspace{0.5cm}

    \textbf{زیرسیستم‌های آزموده‌شده:}
    \begin{itemize}
        \item یکپارچگی \lr{minted} با \lr{Pygments}
        \item کدهای \lr{Python} و \lr{C++}
        \item کد درون‌خطی
        \item کپشن و ارجاع به کد
        \item کامنت فارسی داخل کد
        \item فونت کد \lr{JetBrains Mono}
    \end{itemize}
}

\end{document}